\scanpage{95}
% archive sheet number: 94
% block margin note | boxed upper-right heading
\margin{Impossibilité \unsure{du|de} \unsure{passage} au quotient}
% block text | line 1
$G=GP(N)_{\unsure{\xi|S}}$ opérant sur $P^N_S$\qquad $S=\unsure{\operatorname{Spec}|\mathrm{Sp}}k$

% block text | lines 2-3
$V$ ouvert de $P^N_S$ sur lequel $G$ opère librement, supposons que $V/G$ existe ``effectivement''

% block text | lines 4-12
Alors on montre que le faisceau sur $V/G=W$ déduit de $\underline{L}=\underline{O}(+N+1)=(\underline{\Omega}^N_{P^N_S/S})^{-1}$ par ``descente'' est \uline{ample} sur $W$ [\unsure{utiliser|utilisant} le fait que $G$ n'a pas d'hom. non trivial dans $G_m$, et que \unsure{tt|tout} ouvert affine de $P^N_S$ est complément d'une hypersurface, \deleted{\unsure{donc}} définie par une section d'un $\underline{L}^{\otimes m}$ ($m>0$). À fortiori, $W$ est \uline{affine}. Donc si $G$ opère \inserted{librement} sur $V$ avec un \unsure{graphe} non fermé, $V/G$ n'existe pas.

% block text | lines 13-16, cancelled English block
\deleted{\deleted{\unsure{How to}} \inserted{so as} to free \unsure{arcaïc|archaïc} \unsure{K.-algV|K.-alg} from the \deleted{\unsure{heavy}} \unsure{weight|weights} of \inserted{\deleted{\unsure{unnecessary}} \deleted{\unsure{inadequate}} \unsure{inadequate}} concepts \deleted{conceptual background and technical \unsure{demands|details}} and \deleted{\unsure{new}} tools}
